% docs/assays/rnaseq/methods.tex
\documentclass[11pt]{article}

\usepackage{amsmath, amssymb}
\usepackage{geometry}
\usepackage{hyperref}
\usepackage{booktabs}
\usepackage{enumitem}
\usepackage{longtable}

\geometry{margin=1in}

\title{Supplementary Note: RNA-seq Differential Expression and scRNA Program Loadings to Gene Sets in dig-gene-set-extractors}
\author{}
\date{}

\begin{document}
\maketitle

\section*{Overview and goals}
The \texttt{dig-gene-set-extractors} repository provides converters that map diverse omics data types to:
(i) \emph{gene programs} (scores or weights over genes) and
(ii) compact \emph{gene sets} (typically 100--500 genes) suitable for enrichment and mechanism-level interpretation.

This note specifies the RNA-seq extractors currently implemented as \texttt{rna\_deg} (single contrast), \texttt{rna\_deg\_multi} (batched contrasts from one long table), and \texttt{rna\_sc\_programs} (external scRNA factor-loading ingestion). In contrast to ATAC-seq, RNA-seq is already gene-centric: the primary measured objects are genes (or transcripts), so the central modeling step is not locus-to-gene projection, but rather how to:
(1) define a stable signed gene-level DE score,
(2) transform scores into weights appropriate for downstream use, and
(3) convert scores/weights into one or more gene sets with clear semantics.

\subsection*{Why differential expression as the primary RNA-seq target}
Across bulk and single-cell RNA-seq, the most common user question is contrast-driven:
\begin{itemize}[leftmargin=2em]
\item Which genes are upregulated/downregulated between two conditions (case vs control, treatment vs baseline)?
\item Which genes define a cell type/state relative to others (cluster vs rest)?
\item Which cell-type-specific responses exist (condition A vs B within a cell type)?
\end{itemize}
These contrasts yield signed gene-level statistics that are naturally convertible to gene programs and gene sets.

\subsection*{Replicate-aware note for scRNA-seq}
When biological replicates/donors exist, stable inference is typically achieved by pseudobulk aggregation followed by bulk DE engines (e.g., edgeR, DESeq2, limma-voom), rather than per-cell tests. This note treats the DE results table as input; it does not prescribe a single DE engine. (See Crowell et al.\ for discussion of pseudobulk in scRNA contexts.) \cite{crowell2020muscat}

\section{Mental model: what RNA-seq measures}
RNA-seq measures the abundance of transcripts (or genes) in a sample (bulk) or in individual cells (single-cell).
After normalization and modeling, differential expression analysis produces, for each gene $g$, a signed summary of evidence that expression differs between two conditions.

In a typical two-condition comparison $A$ vs $B$:
\begin{itemize}[leftmargin=2em]
\item An \emph{effect size}, commonly $\ell_g = \log_2(\mathrm{FC}_g)$.
\item A \emph{significance} measure, e.g.\ adjusted p-value $q_g$ (FDR) or p-value $p_g$.
\item Optionally a signed test statistic $t_g$ or $z_g$ (Wald, likelihood ratio, moderated t).
\end{itemize}

The converter maps these into a signed \emph{program score} $s_g$ and then into nonnegative \emph{weights} $w_g$.

\section{How to map CLI arguments to the model in this note}
This note is designed to cross-reference the CLI documentation in \texttt{docs/assays/rnaseq/guide.md}.
Key objects and their corresponding CLI flags are:
\begin{itemize}[leftmargin=2em]
\item Input DE table: \texttt{--deg\_tsv}, with column mapping flags.
\item Signed score definition $s_g$: \texttt{--score\_mode} and column flags (stat/logFC/p).
\item Duplicate gene handling: \texttt{--duplicate\_gene\_policy}.
\item Filtering: \texttt{--exclude\_gene\_regex}, optional \texttt{--gtf} for gene symbol/biotype.
\item Selection of a compact set $S$: \texttt{--select} plus parameters (e.g.\ \texttt{--top\_k}).
\item Weights $w_g$: \texttt{--normalize} (recommended: \texttt{within\_set\_l1}).
\item Gene set export: \texttt{--emit\_gmt}, \texttt{--gmt\_topk\_list}, \texttt{--gmt\_split\_signed}, \texttt{--gmt\_emit\_abs}, \texttt{--gmt\_source}.
\item Naming: \texttt{--signature\_name} (recommended for provenance and interpretability).
\end{itemize}

\section{Formal model}
Let $g \in \mathcal{G}$ index genes. A DE input table provides some subset of:
\[
\ell_g \in \mathbb{R} \quad (\text{log fold-change}), \qquad
q_g \in (0,1] \quad (\text{adjusted p-value}), \qquad
t_g \in \mathbb{R} \quad (\text{signed test statistic}).
\]
We define a signed score:
\begin{equation}
s_g = f(\ell_g, q_g, t_g),
\label{eq:signed_score}
\end{equation}
then define a nonnegative magnitude:
\begin{equation}
m_g = |s_g|,
\label{eq:magnitude}
\end{equation}
select a gene subset $S \subset \mathcal{G}$, and compute weights:
\begin{equation}
w_g =
\begin{cases}
\displaystyle \frac{m_g}{\sum_{h \in S} m_h} & \text{if } g \in S \text{ and using within-set } \ell_1 \text{ normalization}, \\
0 & \text{otherwise.}
\end{cases}
\label{eq:within_set_l1}
\end{equation}
The signed score $s_g$ preserves direction (up vs down), while the weight $w_g$ represents relative importance among selected genes.

\section{Score definitions (the \texttt{score\_mode} axis)}
The DE extractors support the following score definitions. Defaults favor robustness and comparability across sources.

\subsection{S1.0: \texttt{score\_mode=auto} (default)}
Auto mode resolves to:
\begin{itemize}[leftmargin=2em]
\item \texttt{stat}, if a parseable signed statistic is available.
\item otherwise \texttt{logfc\_times\_neglog10p}, if parseable logFC and p/padj are available.
\end{itemize}
If neither condition is met, extraction fails with an actionable column-mapping error.

\subsection{S1.1: \texttt{score\_mode=stat} (preferred when available)}
If a signed test statistic $t_g$ is provided (Wald, moderated t, z, etc.), define:
\[
s_g = t_g.
\]
Rationale: $t_g$ (or $z_g$) already incorporates effect size and variability, and is a natural ranking score.
Common DE engines include such statistics (e.g., DESeq2, edgeR, limma-voom). \cite{love2014deseq2, robinson2010edger, law2014voom}

\subsection{S1.2: \texttt{score\_mode=logfc\_times\_neglog10p} (fallback)}
If $t_g$ is absent but $(\ell_g, q_g)$ are present:
\begin{equation}
s_g = \ell_g \cdot \min\{-\log_{10}(q_g + \epsilon), C\},
\label{eq:logfc_p}
\end{equation}
with a small $\epsilon > 0$ and cap $C$ (e.g., $C=50$) to avoid infinite/overly dominant values when $q_g \approx 0$.
Rationale: incorporates direction and effect size while upweighting strong evidence, but prevents significance from completely dominating.

\subsection{S1.3: \texttt{score\_mode=custom\_column}}
If the user provides a precomputed signed score column, define:
\[
s_g = \texttt{score\_column}(g).
\]
This supports importing external scoring methods (e.g., shrunken logFC, posterior means, Bayesian model outputs).

\subsection{S1.4: duplicate \texttt{gene\_id} aggregation}
If a DE table contains multiple rows for the same \texttt{gene\_id}, the converter aggregates with
\texttt{--duplicate\_gene\_policy}:
\begin{itemize}[leftmargin=2em]
\item \texttt{max\_abs} (default): keep the signed score with largest absolute magnitude.
\item \texttt{sum}: sum signed scores.
\item \texttt{mean}: average signed scores.
\item \texttt{last}: last observed row wins.
\end{itemize}
Duplicate rows are reported in runtime warnings and metadata.

\section{Filtering and gene identity}
\subsection{Gene identifiers and symbols}
The converter operates on a gene identifier string \texttt{gene\_id} and optionally a \texttt{gene\_symbol}.
For GMT export, gene symbols are preferred for human readability, but Ensembl IDs are acceptable if configured.

\subsection{Optional gene annotation via GTF}
If a GTF is provided, annotate genes with:
\begin{itemize}[leftmargin=2em]
\item \texttt{gene\_symbol} (from \texttt{gene\_name})
\item \texttt{gene\_biotype} (from \texttt{gene\_type} or \texttt{gene\_biotype})
\end{itemize}
This enables biotype filtering for GMT output (e.g., protein-coding only).

\subsection{Default artifact filters}
DE signatures often contain technical or broadly expressed artifacts (mitochondrial, ribosomal). A conservative default is:
\begin{itemize}[leftmargin=2em]
\item Exclude mitochondrial genes by symbol prefix: \texttt{\^{ }MT-} (human) or \texttt{\^{ }mt-} (mouse).
\item Exclude ribosomal genes by symbol prefix: \texttt{\^{ }RPS}, \texttt{\^{ }RPL} (human) and optionally \texttt{\^{ }Rps}, \texttt{\^{ }Rpl} (mouse).
\end{itemize}
These filters should be configurable and optional; they can remove real biology in some contexts (e.g., stress responses).
When \texttt{gene\_symbol} is missing, filtering falls back to \texttt{gene\_id}.

\subsection{Auto symbol promotion}
When \texttt{gene\_symbol} is mostly missing, the implementation inspects the \texttt{gene\_id} namespace:
\begin{itemize}[leftmargin=2em]
\item If \texttt{gene\_id} is mostly non-Ensembl-like, \texttt{gene\_symbol} is auto-promoted from \texttt{gene\_id}.
\item If \texttt{gene\_id} is mostly Ensembl-like, no auto-promotion is applied and a warning suggests providing \texttt{--gtf} for symbol mapping.
\end{itemize}

\section{From scores to compact gene sets}
The extractor produces:
\begin{itemize}[leftmargin=2em]
\item A selected program table (\texttt{geneset.tsv}) for reproducibility and downstream weighted analyses.
\item One or more GMT gene sets (\texttt{genesets.gmt}) for enrichment pipelines.
\end{itemize}

\subsection{Selection (the \texttt{select} axis)}
Selection operates on $m_g = |s_g|$ by default to capture the most differential genes regardless of direction.
Common selection strategies:
\begin{itemize}[leftmargin=2em]
\item \texttt{top\_k}: choose the $K$ genes with largest $m_g$.
\item \texttt{threshold}: choose genes with $m_g \ge \tau$.
\item \texttt{quantile}: choose genes above a quantile of $m_g$.
\item \texttt{none}: keep all genes (rarely recommended for GMT).
\end{itemize}

\subsection{Weights and normalization (recommended: within-set L1)}
After selecting $S$, weights are computed from $m_g$ and normalized.
Within-set L1 normalization (Eq.\ \ref{eq:within_set_l1}) yields weights that sum to 1 over selected genes, which is often convenient for downstream methods expecting a distribution.

Importantly, \emph{the signed score $s_g$ is kept separately}, so the representation still supports directionality.

\subsection{Directional gene sets from signed scores}
For enrichment, users typically want separate UP and DOWN sets.
When \texttt{gmt\_split\_signed=true}:
\begin{itemize}[leftmargin=2em]
\item UP set: genes with $s_g > 0$, ranked by $s_g$.
\item DOWN set: genes with $s_g < 0$, ranked by $-s_g$.
\end{itemize}
The GMT export then applies the requested size controls (e.g., \texttt{topk\_list=200} with \texttt{min\_genes=100}, \texttt{max\_genes=500}).

\paragraph{GMT source table.}
By default (\texttt{--gmt\_source=full}), GMT sets are derived from the full ranked score table (\texttt{geneset.full.tsv} when emitted), while \texttt{geneset.tsv} remains the selected compact program used for weighted downstream analyses.
Setting \texttt{--gmt\_source=selected} restricts GMT generation to selected rows only.

\paragraph{Optional directionless set.}
With \texttt{--gmt\_emit\_abs=true}, an additional abs(score)-ranked GMT family is emitted (for users who want a single ``most differential genes'' view independent of sign).

\paragraph{Split-signed edge case.}
If \texttt{--gmt\_split\_signed=true} but no negative scores are present, the converter emits unsplit/positive sets and warns explicitly.

\section{Recommended defaults}
A sensible default configuration for a single DE contrast is:
\begin{itemize}[leftmargin=2em]
\item \texttt{score\_mode=auto} (prefer \texttt{stat} if available, else Eq.\ \ref{eq:logfc_p})
\item \texttt{duplicate\_gene\_policy=max\_abs}
\item \texttt{select=top\_k}, \texttt{top\_k=200}
\item \texttt{normalize=within\_set\_l1}
\item \texttt{emit\_gmt=true}, \texttt{gmt\_topk\_list=200}, \texttt{gmt\_split\_signed=true}, \texttt{gmt\_emit\_abs=false}, \texttt{gmt\_source=full}
\item Conservative artifact filters enabled but user-configurable
\item \texttt{signature\_name=contrast} resolves to the DE filename stem by default (user can override)
\end{itemize}

For batched tables (\texttt{rna\_deg\_multi}), apply the same defaults per comparison, writing one output directory per comparison plus a root manifest.

\section{Practical heuristics and pitfalls}
\begin{itemize}[leftmargin=2em]
\item \textbf{Do not treat per-cell DE as replicate-aware.} Prefer pseudobulk if donors/replicates exist. \cite{crowell2020muscat}
\item \textbf{Column name heterogeneity.} Provide preset mappings for common tools (DESeq2, edgeR, limma, Seurat/Scanpy) or require explicit column flags.
\item \textbf{Gene IDs.} If using Ensembl IDs and you want symbols in GMT, provide a GTF or a symbol mapping; otherwise set \texttt{gmt\_prefer\_symbol=false} and \texttt{gmt\_require\_symbol=false}.
\item \textbf{Extreme p-values.} Cap $-\log_{10}(q_g)$ and add $\epsilon$.
\end{itemize}

\section{scRNA program loadings (external factorization outputs)}
The \texttt{rna\_sc\_programs} converter ingests precomputed gene loadings from external methods (for example cNMF and scHPF outputs). It does \emph{not} fit factor models itself.

Let $l_{pg}$ denote loading for program (factor) $p$ and gene $g$. The converter defines transformed program score $s_{pg}$ by \texttt{score\_transform}:
\[
s_{pg} =
\begin{cases}
l_{pg} & \texttt{signed}\\
|l_{pg}| & \texttt{abs}\\
\max(l_{pg}, 0) & \texttt{positive}\\
\max(-l_{pg}, 0) & \texttt{negative}
\end{cases}
\]
Default is \texttt{positive}, yielding nonnegative program membership-like scores and avoiding sign ambiguity for first-pass interpretation.

\subsection*{Per-program selection and weights}
Selection is performed independently for each program:
\[
S_p = \mathrm{Select}\left(\{g, b_{pg}\}_g\right),
\]
where
\[
b_{pg} =
\begin{cases}
|s_{pg}| & \text{if } \texttt{score\_transform=signed}\\
s_{pg} & \text{otherwise}.
\end{cases}
\]
The default \texttt{top\_k=200} is applied per program, and weights use within-program $\ell_1$ normalization:
\[
w_{pg} = \frac{b_{pg}}{\sum_{h \in S_p} b_{ph}}.
\]

\subsection*{Grouped outputs}
Each program is emitted as its own directory under a grouped root:
\begin{itemize}[leftmargin=2em]
\item \texttt{manifest.tsv}
\item \texttt{program=<ID>/geneset.tsv}
\item \texttt{program=<ID>/geneset.meta.json}
\item optional \texttt{program=<ID>/geneset.full.tsv}
\item optional \texttt{program=<ID>/genesets.gmt}
\end{itemize}
An optional combined \texttt{genesets.gmt} is written at the root.

\subsection*{Signed loadings and GMT directionality}
When loadings are signed and the user requests \texttt{gmt\_split\_signed=true}, GMT emits \texttt{\_\_pos} and \texttt{\_\_neg} sets per program.
Under default \texttt{score\_transform=positive}, negatives are dropped and only positive sets are emitted. The converter warns when a substantial fraction of raw loadings are negative under this default.

\subsection*{Practical heuristics for fitting programs on large atlases}
For very large scRNA maps, upstream factorization choices usually dominate downstream gene-set quality.
\begin{itemize}[leftmargin=2em]
\item \textbf{Fit within cell types or broad compartments.} A single global factorization across heterogeneous cell states often produces diffuse programs.
\item \textbf{Downsample before fitting.} Cap total cells and cap cells per donor (and ideally per cell type) to keep both runtime and donor influence controlled.
\item \textbf{Control donor imbalance.} If one donor contributes most cells, factors can become donor-driven rather than biology-driven.
\item \textbf{Keep gene filtering consistent.} Use the same mitochondrial/ribosomal policy upstream and during conversion.
\item \textbf{Choose moderate factor count first.} Start with conservative $K$, then increase only if resulting programs remain interpretable.
\item \textbf{Respect cNMF input constraints.} Remove zero-total cells and zero-total genes before factorization; cNMF warnings on these inputs usually indicate unstable factorization setup.
\end{itemize}
To support practical workflows, this repo provides \texttt{geneset-extractors workflows scrna\_cnmf\_prepare}, which performs deterministic downsampling, subset splitting, and zero-total filtering, and writes per-subset cNMF run scripts. Separately, \texttt{rna\_sc\_programs} emits non-fatal warnings when very large numbers of programs or parsed loading values are detected during ingestion.

\subsection*{Auto-selection of factor count ($K$)}
For cNMF workflows, $K$ is selected from the cNMF stability/error diagnostics rather than fixed globally.
The repo workflow command \texttt{geneset-extractors workflows cnmf\_select\_k} reads the cNMF
\texttt{k\_selection\_stats} output and applies a reproducible default heuristic:
\begin{itemize}[leftmargin=2em]
\item Strategy: \texttt{largest\_stable}
\item Candidate rule: stability $\ge 0.95 \times \max(\text{stability})$
\item Choice: largest $K$ among candidates
\item Fallback: max-stability $K$ when no candidate passes thresholds
\end{itemize}
This preserves practical flexibility (no universal single criterion) while keeping the default automated and reproducible.
Users can override with explicit fixed $K$ or stricter local-maximum constraints.

\begin{thebibliography}{9}
\bibitem{love2014deseq2}
Love MI, Huber W, Anders S.
Moderated estimation of fold change and dispersion for RNA-seq data with DESeq2.
Genome Biology (2014).

\bibitem{robinson2010edger}
Robinson MD, McCarthy DJ, Smyth GK.
edgeR: a Bioconductor package for differential expression analysis of digital gene expression data.
Bioinformatics (2010).

\bibitem{law2014voom}
Law CW, Chen Y, Shi W, Smyth GK.
voom: precision weights unlock linear model analysis tools for RNA-seq read counts.
Genome Biology (2014).

\bibitem{crowell2020muscat}
Crowell HL et al.
muscat detects subpopulation-specific state transitions from multi-sample multi-condition single-cell transcriptomics data.
Nature Communications (2020).
\end{thebibliography}

\end{document}
